% Example of a vector figure, meant to be included with \input{} inside a figure environment
% Such fragments are typically produced by a script, so that data and rendering stay in sync
% The tikzpicture environment is patched by the template to never overflow \linewidth

\begin{tikzpicture}

    \begin{axis}
        [
            height      = 5cm,
            width       = \linewidth,
            xlabel      = {Number of training epochs},
            ylabel      = {Placeholder metric},
            xmin        = 0, xmax = 20,
            ymin        = 0, ymax = 1,
            grid        = major,
            grid style  = {separatorColor},
            legend pos  = south east,
            legend cell align = left,
        ]

        \addplot[titlesColor, line width=1.5pt, mark=*, mark size=1.5pt] coordinates
            {
                (0, 0.10) (2, 0.38) (4, 0.55) (6, 0.66) (8, 0.73)
                (10, 0.78) (12, 0.82) (14, 0.84) (16, 0.86) (18, 0.87) (20, 0.87)
            };
        \addlegendentry{Placeholder model}

        \addplot[linksColor, line width=1.5pt, dashed, mark=square*, mark size=1.5pt] coordinates
            {
                (0, 0.10) (2, 0.30) (4, 0.44) (6, 0.53) (8, 0.59)
                (10, 0.63) (12, 0.66) (14, 0.68) (16, 0.69) (18, 0.70) (20, 0.70)
            };
        \addlegendentry{Placeholder baseline}

    \end{axis}

\end{tikzpicture}
